\subsection*{从一组明确的初始数据开始}

Alice 使用账户 1，初始余额为 10000 分，购买、订单和扣款记录为空。商品 \code{paper-001} 售价 3000 分，币种为 CNY。

购买接口保留第一章的规则。\code{POST /purchases} 向购买入口提交一次购买请求。每次受理，服务生成一个新的 \code{request\_id}，用它标记本次请求尝试；进入购买处理后再生成新的 \code{purchase\_id}，用它标记将要保存的购买记录。

购买成功时，条件 \code{UPDATE} 在余额至少达到商品价格时扣减。服务随后在同一 PostgreSQL 数据库、同一连接、同一显式事务中保存购买、订单和扣款流水，最后用 COMMIT 提交。这组写入仍遵守第一章的原子性：共同生效或共同撤回。

购买最终为 \code{succeeded}（成功）或 \code{insufficient\_funds}（余额不足）。这两项是本例的终态：一旦提交保存，这笔购买就以相应结果结束；以后的查询沿用这个结果。\code{pending} 表示处理中，只在事务内部使用。服务收到提交成功确认后，再向客户端返回业务结果。

把第一次 HTTP 受理简称为 R1，对应的购买、订单和流水分别简称为 P1、O1、D1。这些简称表示实际各自独立的编号。第二次受理叫 R2，业务编号相应为 P2、O2、D2。先让请求串行处理：R1 的数据库事务结束以后，客户端才发出 R2。这样的次序已经足以产生两次扣款。

客户端发送的仍是同一份请求体，JSON 中的商品编号与第一次完全相同：

\begin{BookCode}
POST /purchases HTTP/1.1
Authorization: Bearer <access_token>
Content-Type: application/json

{"product_id": "paper-001"}
\end{BookCode}

报文中的 \code{Authorization} 携带访问凭据，服务验证后据此取得 Alice 的账户；\code{Content-Type} 说明请求体采用 JSON 格式。响应也有说明状态和内容格式等信息的报文头，称为响应头，具体业务数据放在其后的响应体中。

服务按身份取得账户 1，按商品目录取得金额 3000。它创建 P1，执行条件扣款，把余额从 10000 改为 7000，写入 O1 和 D1，再把 P1 更新为 \code{succeeded}。这些写入目前属于尚未结束的事务。最终提交成功以后，它们才共同成为后续查询能够读取的业务历史。

\subsection*{故障只放在提交以后}

这次故障的注入位置很具体：\textbf{数据库已经提交，服务也已收到提交成功确认，但客户端尚未收到完整的 HTTP 成功响应。}先在这里暂停响应路径，让客户端在等待期限内收不到响应，也没有及时收到连接关闭通知。

服务立即退出或主动关闭连接时，客户端可能很快收到 EOF 或连接重置。EOF 是 End of File，在这里表示读取到了连接数据的末尾；连接重置表示连接被异常中止。它们都可能使读取立刻结束。这样的观察也意味着客户端尚未取得业务结果，但表现与等待期限耗尽不同。这里让连接保持到客户端超时，再丢弃本次响应，使章首的故障有一条明确的执行路径。

语言无关的控制流程如下。故障点只用于分析这次历史，正常处理仍然在提交确认后返回结果。

\begin{BookCode}
受理 R1，准备 P1、O1、D1
开始数据库事务
    保存购买输入
    条件更新账户余额
    保存订单、扣款流水和购买终态
提交事务
收到数据库提交成功确认
    ← 在这里暂停并最终丢弃本次响应
向客户端返回 P1、O1、D1
\end{BookCode}

若服务在保存订单之后、提交之前退出，数据库结束未提交事务时会撤回相关变化。这里的事务已经完成提交，故障作用于后面的响应传递。服务即使随后退出，内存中的待发送内容消失，数据库仍按它的持久性保证保留已经确认的购买。

\begin{figure}[htbp]
\centering
\begin{tikzpicture}[x=1cm,y=1cm]
\foreach \x in {0,5,10} {\draw[Ink!35,dashed] (\x,-.4)--(\x,-6.4);}
\node[booknode,text width=2cm] at (0,0) {客户端};
\node[booknode,text width=2cm] at (5,0) {购买服务};
\node[booknode,text width=2cm] at (10,0) {数据库};
\draw[bookarrow] (0,-.9)--node[above,font=\footnotesize]{R1：购买 paper-001}(5,-.9);
\draw[bookarrow] (5,-1.8)--node[above,font=\footnotesize]{购买事务与 COMMIT}(10,-1.8);
\node[booknode,text width=3.5cm] at (10,-2.6) {P1、O1、D1\\已提交，余额 7000};
\draw[bookarrow] (10,-3.4)--node[above,font=\footnotesize]{提交成功确认}(5,-3.4);
\node[bookresult,text width=3.6cm] at (5,-4.3) {暂停响应路径\\客户端尚未收到结果};
\draw[bookalert,dashed] (5,-5.2)--node[below,font=\footnotesize]{本次响应丢失}(1.4,-5.2);
\draw[Accent,line width=1.2pt] (1.1,-5)--(1.5,-5.4) (1.1,-5.4)--(1.5,-5);
\node[bookresult,text width=2.4cm] at (0,-6.1) {等待期限耗尽};
\node[booknote,text width=3.2cm] at (10,-5.4) {响应中断发生在\\业务记录提交以后};
\end{tikzpicture}
\caption{提交已经完成，结果传递中断。客户端超时以后，数据库中仍有一次成功购买。}
\label{fig:ch02-commit-gap}
\end{figure}


余额为 7000，P1 的状态为 \code{succeeded}，O1 为 \code{paid}（已付款），D1 的扣除额为 3000、当次扣后余额为 7000。客户端经历的则是发送请求、持续等待，最后到达本地截止时间。它尚未拿到能够解析并确认的业务结果。

此时从独立数据库连接读取，可以看到余额 7000，购买、订单、流水各一条。这个读取不依赖客户端收到响应：事务已经提交，暂停的是后面的结果传递。本章假定数据库保留已经确认的提交；同步提交与存储恢复的具体前提将在第 2.3 节说明。

\subsection*{第二次点击为什么又扣了钱}

Alice 看到超时后再次点击购买。客户端建立新的 HTTP 请求，发送完全相同的 JSON。服务按现有规则给这次受理分配 R2，创建新购买 P2，随后执行同样的事务。

账户当前余额为 7000，仍然足以支付 3000。条件 \code{UPDATE} 成功，余额变为 4000。P2、O2、D2 共同提交。第二次响应正常到达，客户端取得 P2、O2、D2，显示购买成功。第一次的 P1 仍留在数据库中，R1 的连接中断发生在这笔购买提交以后。

\begin{center}
\begin{minipage}{\linewidth}
\centering\small\captionsetup{hypcap=false}
\captionof{table}{两次请求留下的记录}
\begin{tabularx}{\linewidth}{p{.24\linewidth}XX}
\toprule
\textbf{观察位置} & \textbf{R1 提交并丢失响应后} & \textbf{R2 正常完成后} \\
\midrule
账户余额 & 7000 & 4000 \\
成功购买 & P1 & P1、P2 \\
已付款订单 & O1 & O1、O2 \\
扣款流水 & D1：3000 & D1：3000；D2：3000 \\
客户端收到的成功结果 & 无 & P2、O2、D2 \\
\bottomrule
\end{tabularx}
\end{minipage}
\end{center}

账本仍然满足第一章的核对规则：10000 减去两条流水合计的 6000，得到 4000。每笔购买最多对应一张订单，每张订单最多对应一条扣款流水，外键也全部成立。外键要求记录中引用的购买或订单确实存在，因而两组完整记录都能通过检查。逐笔检查购买时，两组记录都符合数据库约束。Alice 只打算购买一次，账户却为此减少了 6000 分。

\begin{figure}[htbp]
\centering
\begin{tikzpicture}[x=1cm,y=1cm]
\node[booknode,text width=2.1cm] (r1) at (0,0) {R1\\第一次请求};
\node[booknode,text width=3.2cm] (p1) at (4.4,0) {P1 $\to$ O1 $\to$ D1\\扣款 3000};
\node[booknode,text width=2.5cm] (b1) at (9.6,0) {余额\\10000 $\to$ 7000};
\draw[bookarrow] (r1)--(p1);
\draw[bookarrow] (p1)--(b1);
\node[bookresult,text width=2.1cm] (r2) at (0,-2) {R2\\超时后重试};
\node[bookresult,text width=3.2cm] (p2) at (4.4,-2) {P2 $\to$ O2 $\to$ D2\\再扣 3000};
\node[bookresult,text width=2.5cm] (b2) at (9.6,-2) {余额\\7000 $\to$ 4000};
\draw[bookalert] (r2)--(p2);
\draw[bookalert] (p2)--(b2);
\draw[bookarrow] (b1.south)--(b2.north);
\draw[bookalert,dashed] (r1.south)--node[right,font=\scriptsize]{未取得结果后再发}(r2.north);
\node[booknote,text width=4.4cm] at (4.4,-3.4) {每个购买编号各对应一张订单\\P1、P2 都满足现有唯一约束};
\node[booknote,text width=3.3cm] at (9.6,-3.4) {流水合计 6000\\余额与账本仍相符};
\end{tikzpicture}
\caption{两个请求建立两组独立记录。数据库约束检查各组内部关系，Alice 的一次购买意图尚未作为跨请求关系保存。}
\label{fig:ch02-two-purchases}
\end{figure}


P1 在服务受理第一次请求后产生，客户端原本要从响应取得它。响应丢失以后，客户端手里没有 P1，第二次请求又没有携带可以指向 P1 的信息。服务看到两个通过授权、参数合法的 POST，依照现有规则建立了两笔购买。

同样的 JSON 也可能来自两次独立购买。Alice 收到第一次成功响应，决定再买一包相纸，随后再次发送 \code{paper-001}。服务若把所有相同 JSON 都合并，这两次购买就只能成立一次。将两次点击间隔限制在十秒内，也会遇到用户十秒内确实购买两包、旧请求十秒以后才重试的情况。

\begin{figure}[htbp]
\centering
\begin{tikzpicture}[x=1cm,y=1cm]
\node[booknode,text width=9cm] at (5.6,0) {两次请求体相同：\texttt{\{"product\_id":"paper-001"\}}};
\node[font=\small\sffamily\bfseries] at (2,-1) {响应丢失后的再次尝试};
\node[font=\small\sffamily\bfseries] at (9.2,-1) {买完一包，再买一包};
\node[booknode,text width=3cm] (u) at (2,-1.9) {一个购买意图 $u$};
\node[booknode,text width=1.4cm] (a) at (.5,-3.2) {R1};
\node[booknode,text width=1.4cm] (b) at (3.5,-3.2) {R2};
\draw[bookarrow] (u)--(a);
\draw[bookarrow] (u)--(b);
\node[booknode,text width=2.6cm] (u1) at (7.5,-1.9) {购买意图 $u_1$};
\node[booknode,text width=2.6cm] (u2) at (10.9,-1.9) {购买意图 $u_2$};
\node[booknode,text width=1.4cm] (c) at (7.5,-3.2) {R1};
\node[booknode,text width=1.4cm] (d) at (10.9,-3.2) {R2};
\draw[bookarrow] (u1)--(c);
\draw[bookarrow] (u2)--(d);
\node[bookresult,text width=4.4cm] at (2,-4.6) {预期：一笔购买\\累计扣除 3000};
\node[booknode,text width=4.4cm] at (9.2,-4.6) {预期：两笔购买\\累计扣除 6000};
\end{tikzpicture}
\caption{相同参数既可能属于同一操作的重试，也可能属于两次独立购买。区分两者需要请求体以外的操作关系。}
\label{fig:ch02-same-input}
\end{figure}


重复扣款的判断因此需要结合 Alice 发起购买的目的：她授权购买一包相纸，第二次发送用于取得或完成这次购买。另一种情况下，她明确决定再买一包，两笔购买就都应成立。服务需要有信息区分这两种关系，商品、数量和点击间隔都不足以独自完成这个判断。

\subsection*{为什么调整响应与提交的顺序没有解决问题}

服务在提交确认以后返回成功，中间的间隔似乎可以再缩短一些。缩短间隔减少了这一阶段的暴露时间；若要进一步计算故障发生概率，还须给出故障到达的随机模型。数据库保存记录以后，结果仍要经过服务发送、网络传输、客户端接收和解释，这些动作可以分别中断。

若先发送成功响应，再提交数据库，客户端可能已经显示成功，事务却在随后遇到错误并回滚。这时客户收到的成功缺少已提交购买的支撑。按当前接口规则，应当先有已提交的业务事实，再报告成功。

若一直持有数据库事务，等客户端确认收到结果以后才提交，客户端确认本身也要经过网络。确认丢失以后，服务又面对一次等待：客户究竟没有收到成功，还是已经收到、只是确认没送回来？此前已经向客户报告购买完成，随后回滚，还会使已报告的结果失去对应记录。

\begin{figure}[htbp]
\centering
\begin{tikzpicture}[x=1cm,y=1cm]
\node[booknote,anchor=west] at (-1.1,.7) {先提交，再报告成功};
\node[booknode,text width=2.6cm] (a) at (.4,0) {数据库提交成功};
\node[bookresult,text width=2.6cm] (b) at (4.9,0) {响应传递中断};
\node[booknode,text width=3cm] (c) at (9.8,0) {已扣款\\客户端结果未知};
\draw[bookarrow] (a)--(b);
\draw[bookalert] (b)--(c);
\node[booknote,anchor=west] at (-1.1,-1.3) {先报告成功，再提交};
\node[booknode,text width=2.6cm] (d) at (.4,-2) {客户端收到成功};
\node[bookresult,text width=2.6cm] (e) at (4.9,-2) {事务随后回滚};
\node[bookresult,text width=3cm] (f) at (9.8,-2) {报告了成功\\数据库却无购买};
\draw[bookarrow] (d)--(e);
\draw[bookalert] (e)--(f);
\node[booknote,anchor=west] at (-1.1,-3.3) {等待客户端再确认};
\node[booknode,text width=2.6cm] (g) at (.4,-4) {先发送结果};
\node[bookresult,text width=2.6cm] (h) at (4.9,-4) {客户端确认丢失};
\node[booknode,text width=3cm] (i) at (9.8,-4) {服务仍需决定\\怎样继续事务};
\draw[bookarrow] (g)--(h);
\draw[bookalert] (h)--(i);
\end{tikzpicture}
\caption{调整先后次序会改变故障后果。成功响应需要已提交事实支撑，结果丢失则需要后续恢复路径。}
\label{fig:ch02-ordering}
\end{figure}


即使服务写出全部响应字节的调用已经返回，客户端也可能尚未读取完整响应，或者读入内存后尚未保存购买编号就退出。写入调用通常只说明本地网络栈接受了这些字节。远端应用取得并保存了什么，需要由相应协议和后续证据确认。

购买服务需要规定哪些事实必须持久保存、恢复时从哪里查询，以及重复请求如何处理。
\FloatBarrier
